\section{対象と結果の範囲}\label{sec:conventions}

係数体を \(k=\C\) とする。代数は連結局所有限な
\(\mathbb N\)-次数付き \(k\)-代数
\(A=\bigoplus_{j\ge0}A_j\)、\(A_0=k\) とする。
\(A_+=\bigoplus_{j>0}A_j\)、\(k=A/A_+\) と書き、
shift は \(M(a)_j=M_{j+a}\) とする。
\(H_M(t)=\sum_j(\dim_kM_j)t^j\) は、各次数有限次元で次数下有界な
次数付きベクトル空間 \(M\) の Hilbert 級数である。

AS-regular とは、大域次元 \(d\) と GK 次元が有限であり、
左右の AS-Gorenstein 条件を満たすことをいう。
右加群の規約では、その条件は
\[
\underline{\Ext}^{i}_A(k,A)\cong
\begin{cases}k(\ell)&i=d,\\0&i\ne d\end{cases}
\]
である。\(\underline{\Ext}\) はすべての内部次数を含む次数付き Ext を表す。
\(\ell=\ASind A\) を AS index と呼ぶ。
本文の構造定理では、両側 noetherian 性を仮定する。

\emph{最小生成次数が \((1^n,2)\)} とは、
\(A_+/A_+^2\) が次数 \(1\) に \(n\) 次元、次数 \(2\) に \(1\) 次元を持ち、
他の次数では零であることをいう。
特に、次数 \(2\) の生成元は次数 \(1\) の元の積で消去できない。
\[
V=A_1,\qquad B=k\langle A_1\rangle\subseteq A
\]
と書く。\(B\) はこの定義から \(A\) に内在的に定まる。
自由テンソル代数を \(T(V)=\bigoplus_{j\ge0}V^{\otimes_k j}\)、
その次数 \(j\) の部分を \(T_j(V)=V^{\otimes_k j}\) と書く。

次数 \(2\) の Ore 拡大の規約は
\[
A=B[w;\sigma,\delta],\qquad
wb=\sigma(b)w+\delta(b),\qquad |w|=2,
\]
である。ここで \(\sigma\in\Aut_{\mathrm{gr}}B\)、
\(\delta(B_j)\subseteq B_{j+2}\) であり、
\[
\delta(ab)=\sigma(a)\delta(b)+\delta(a)b.
\]
この拡大は左 \(B\)-加群として \(\bigoplus_{j\ge0}Bw^j\) であり、
\(\sigma\) が可逆なので右 \(B\)-加群としても
\(\bigoplus_{j\ge0}w^jB\) である。

主結果である定理~\ref{thm:ore} は、
\(\gldim A=n+1\)、\(\ASind A=n+2\) の下で、この Ore 表示と
\(B\) の \(n\) 次元 quadratic Koszul AS-regular 性を与える。
\(n=2,3\) では、必要な AS index を大域次元と最小生成次数から決定できる
（第~\ref{sec:lowdim} 節）。従って、とくに 4 次元 \((1,1,1,2)\) 型は、
3 次元 quadratic AS-regular 代数とその Ore データに帰着する。
その記述の同型関係も明記する。

一般の \(n\) について、最小生成次数と \(\gldim A=n+1\) のみから
\(\ASind A=n+2\) を導く主張は、このノートでは証明しない。
第~\ref{sec:counterexample} 節の反例は大域次元 \(n+2\) を持つため、
大域次元 \(n+1\) に限定した主張への反例ではない。
